\documentclass{article}
% NeurIPS 2026 workshop submission (single blind).
\usepackage[dblblindworkshop]{neurips_2026}
% TODO: fill in the workshop name before submitting.
\workshoptitle{PUDM}

\usepackage[utf8]{inputenc}
\usepackage[T1]{fontenc}
\usepackage{hyperref}
\usepackage{url}
\usepackage{booktabs}
\usepackage{amsfonts}
\usepackage{amsmath}
\usepackage{amssymb}
\usepackage{nicefrac}
\usepackage{microtype}
\usepackage{xcolor}
\usepackage{graphicx}
\usepackage{amsthm}
\usepackage{enumitem}

\newtheorem{theorem}{Theorem}
\newtheorem{proposition}{Proposition}
\newtheorem{assumption}{Assumption}
\newtheorem{lemma}{Lemma}
\newtheorem{corollary}{Corollary}

\newcommand{\method}{CP-IRL}
\newcommand{\Dnorm}[1]{\|#1\|_{\mathcal D}}

\title{Calibrating What Behavior Identifies:\\
Conformal Reward Sets for Inverse Reinforcement Learning}

% Double-blind: identity intentionally omitted from the source.
\author{}

\begin{document}

\maketitle

\begin{abstract}
Conformal inverse optimization (CIO) calibrates an uncertainty set around an
estimated objective, then optimizes robustly against it. We ask which parts of
that recipe survive the move to inverse reinforcement learning (IRL). The
conformal calibration transfers untouched; two other components break, both
because behavior identifies a reward only up to transformations that leave every
policy's value unchanged. CIO scores a demonstration by an angle between reward
coordinate vectors, and a feature reparameterization that changes no reward
value and no demonstrator's policy still moves the calibrated radius by
$24^\circ$; we replace the angle by distance to a Bellman-feasible reward cone
in the span seminorm of the occupancy polytope, which is invariant by
construction and still a linear program. The calibrated event is also weaker
than it looks: it certifies that \emph{some} reward in the set explains the new
demonstrator ($0.82$ measured, at target $0.80$), not that the demonstrator's
own reward lies in it ($0.44$). Closing that gap takes identification rather
than more data---a certified bound $\eta$ on inverse-fiber width---giving the
containment radius $R=\min\{2,2q_\gamma+\eta\}$ and a sharp threshold $R\ge1$
past which the safety guarantee provably carries no decision content.
\end{abstract}

\section{Introduction}

IRL turns behavior into a reward estimate and then optimizes that estimate,
discarding the uncertainty left by the first step exactly where it matters:
different rewards consistent with the same demonstrations prescribe different
policies. Conformal inverse optimization (CIO) is an appealing remedy for
one-shot inverse problems \citep{lin2024conformal}. Observe a population of
decision makers, fit a point estimate on part of it, and on the rest measure how
far that estimate sits from the set of objectives that would rationalize each
held-out decision; the conformal quantile of those distances defines an
uncertainty set with finite-sample coverage, and robust optimization over the
set replaces the point prescription. The recipe is modular and looks portable.
We carried it into IRL and report what happened, component by component
(Table~\ref{tab:recipe}).

\begin{table}[t]
\centering
\small
\caption{The CIO recipe, component by component, in IRL. Bold entries mark what
had to change; the rest, including the conformal quantile itself, are inherited
and used as such.}
\label{tab:recipe}
\begin{tabular}{@{}p{0.27\linewidth}p{0.40\linewidth}p{0.26\linewidth}@{}}
\toprule
CIO ingredient & Counterpart here & Status \\
\midrule
Point estimate; inverse-feasible set of a decision & any IRL fit; the Bellman
cone $K(\pi)$, linear in $\theta$ after resolvent substitution & inherited \\
Score: angle to the estimate & span-seminorm distance to $K(\pi)$; invariant,
still an LP & \textbf{redesigned} (\S\ref{sec:method}, E1) \\
Split-conformal quantile & unchanged & inherited, textbook \\
Calibrated event $\Rightarrow$ parameter set & needs an inverse-fiber width
$\eta$ that behavior cannot supply & \textbf{new assumption}
(\S\ref{sec:bridge}, E2) \\
Robust decision & reference-relative robust MDP, a single LP; absolute version
is unbounded & \textbf{redesigned} (\S\ref{sec:bridge}) \\
\bottomrule
\end{tabular}
\end{table}

One component transfers for free, and it is not our contribution: once
per-demonstrator scores are defined, the rest of the calibration is split
conformal prediction---rank the scores, take the $\lceil\gamma(N+1)\rceil$-th,
invoke exchangeability \citep{vovk2005algorithmic,lin2024conformal}. Sequential
decision making changes no step of that argument. The work is in the components
on either side, and both break for the same reason.

\paragraph{Behavior does not fix reward coordinates.} Potential shaping, rewards
on unreachable states, and other additive directions change a reward vector
while adding the same constant to every policy's value, so an angle between
reward coordinate vectors spends its radius on distinctions no demonstration can
reveal. This is not hypothetical: an invertible feature reparameterization
leaving every scalar reward and every demonstrator's policy pointwise unchanged
still moves the calibrated angular radius by $24^\circ$ (E1), so two
mathematically identical problems get two different uncertainty sets. The fix is
not a better estimator but an answer to the prior question---what does behavior
distinguish?

\paragraph{Intersection is not containment.} Split conformal certifies that
\emph{some} reward in the calibrated set makes the new demonstrator optimal; a
robust prescription needs the demonstrator's \emph{own} reward to lie in the
set. In one-shot inverse optimization the two nearly coincide, but in IRL a
single policy is rationalized by a whole cone: at target $0.8$ we measure the
first event at $0.82$ and the second at $0.44$ (E2). Calibration cannot close
that gap, because the missing quantity is identification, not statistics.

We therefore rebuild the score in decision-space geometry
(\S\ref{sec:method}) and add an explicit bridge from intersection to containment
(\S\ref{sec:bridge}), with a sharp boundary at which its own safety promise
becomes vacuous. We do \emph{not} claim that \method{} lowers true-reward regret
against point-estimate IRL; the claim is narrower and, for anyone considering
the same transplant, more useful: half of CIO's pipeline moves to IRL unchanged,
two components must be rebuilt, and one thing cannot be supplied by behavioral
data at all.

\paragraph{Related work.} CIO calibrates uncertainty for a latent objective
\citep{lin2024conformal} using split conformal prediction
\citep{vovk2005algorithmic}. Two lines are directly relevant: partial
identifiability of reward sets in IRL \citep{skalse2023invariance,
kim2024unified} and reward-comparison metrics, of which STARC
\citep{skalse2024starc} subsumes EPIC \citep{gleave2021epic}. Our occupancy span
is the canonical seminorm behind one STARC construction, so we claim no new
reward metric: the contribution is calibrating Bellman cones \emph{in} that
geometry, plus the fiber-width bridge. Appendix~\ref{app:related} covers the
rest.

\section{An invariant score with exact intersection coverage}
\label{sec:method}

A tabular MDP $M=(S,A,P,\rho_0,\beta,\phi)$ is shared and known, with rewards
$r_\theta(s,a)=\phi(s,a)^\top\theta$; demonstrator $k$ draws
$\hat\theta_k\sim P_\theta$, acts optimally under it, and we observe its
deterministic policy $\pi_k$ exactly. Let $M(P)$ be the polytope of normalized
discounted occupancy measures, $\Phi$ the feature matrix,
$V_\theta(\mu):=\theta^\top\Phi^\top\mu/(1-\beta)$, and let $K(\pi):=\{\theta:
Q^\pi_\theta(s,a)\le Q^\pi_\theta(s,\pi(s))\ \forall s,a\}$ be the unbounded
Bellman cone of rewards making $\pi$ optimal; fixing $\pi$ makes it exactly
$|S|(|A|-1)$ halfspaces (Proposition~\ref{prop:resolvent}).

\begin{assumption}[Exchangeable calibration unit]
\label{ass:unit}
The pairs $(\hat\theta_k,\pi_k)$, $k=1,\dots,N,N{+}1$, are i.i.d.\ (hence
exchangeable), and $M$ (in particular $P$) is fixed across all demonstrators and
the deployment instance.
\end{assumption}

One demonstrator is thus one calibration point, playing the role of CIO's
(decision, context) pair.

\paragraph{The occupancy-difference span.} For $v\in\mathbb R^d$ let
\[
\Dnorm{v}:=\max_{\mu\in M(P)}V_v(\mu)-\min_{\mu\in M(P)}V_v(\mu)
= V^*_v(\rho_0)+V^*_{-v}(\rho_0),
\]
two value-iteration solves. Read $\Dnorm{v}$ as: how much can the reward
direction $v$ change any policy's value? Directions no policy can exploit score
zero.

\begin{proposition}[Properties of $\Dnorm{\cdot}$]
\label{prop:da-seminorm}
$\Dnorm{\cdot}$ is a seminorm with $\Dnorm{v}=0$ iff $v\in\mathcal A(P):=\{a:
a^\top\Phi^\top(\mu-\nu)=0\ \forall\mu,\nu\in M(P)\}$, the \emph{additively
null} subspace, which contains the potential-shaping subspace
$S_{\mathrm{pot}}$ and can be strictly larger. For any $\theta,\theta'$ with
$\mu_{\theta'}$ optimal for $\theta'$,
$\mathrm{Reg}_\theta(\pi_{\theta'}):=V_\theta(\mu^*_\theta)-V_\theta(\mu_{\theta'})
\le\Dnorm{\theta-\theta'}$. Finally $\Dnorm{\cdot}$ is exactly invariant under
$\phi\mapsto A\phi,\ \theta\mapsto A^{-\top}\theta$ for invertible $A$.
\end{proposition}

The regret constant is $1$---no Cauchy--Schwarz slack, no feature-norm Lipschitz
factor---because $\Dnorm{\cdot}$ is defined as the very value spread it must
bound. Behavior also cannot identify positive reward scale, which we remove by
normalizing rather than by building quotient coordinates: writing
$\theta^\circ:=\theta/\Dnorm{\theta}$, the map
$d_\circ(\theta,\theta'):=\Dnorm{\theta^\circ-\theta'^\circ}$ is a pseudometric
on $[0,2]$ that vanishes exactly when $\theta'=c\theta+a$ with $c>0$,
$a\in\mathcal A(P)$, and still bounds normalized regret
(Proposition~\ref{prop:projective}). So $d_\circ$ compares rays of policy-value
differences.

\paragraph{Score and calibrated set.} With $\bar\theta$ an IRL point estimate
fit on a disjoint training split, define
\[
s_k:=\min_{\theta\in K(\pi_k)}\Dnorm{\theta-\bar\theta^\circ},
\qquad B_q:=\{\theta:\Dnorm{\theta-\bar\theta^\circ}\le q\}.
\]
Since $\Dnorm{\cdot}$ has a linear-program epigraph, $s_k$ is a single LP in
$d+2|S|$ variables, with no second-order cone anywhere; the Euclidean ball CIO
uses instead would make it an SOCP and destroy reparameterization invariance.
Because $0\in K(\pi_k)$ we have $s_k\in[0,1]$, and both score and quantile are
unchanged under $\bar\theta\mapsto c\bar\theta+a$ for $c>0$,
$a\in\mathcal A(P)$.

\begin{theorem}[Decision-aware coverage, exact equivalence]
\label{thm:da-coverage}
Calibrate $q_\gamma$ as the $\tau$-th smallest $s_k$ among $N_{\mathrm{val}}$
demonstrators, $\tau=\lceil\gamma(N_{\mathrm{val}}+1)\rceil$ (sentinel $+\infty$
if $\tau=N_{\mathrm{val}}+1$). For a fresh i.i.d.\ demonstrator,
$s_{\mathrm{new}}\le q_\gamma\iff K(\pi_{\mathrm{new}})\cap
B_{q_\gamma}\ne\emptyset$, and
$\Pr[K(\pi_{\mathrm{new}})\cap B_{q_\gamma}\ne\emptyset]\ge\gamma$.
\end{theorem}

The probability statement is the standard split-conformal rank argument
\citep{vovk2005algorithmic} applied verbatim. What is specific here is the
equivalence to its left, immediate from $s_k$ being a minimum: coverage of the
score \emph{is} intersection coverage of the set, with no lossy conversion in
between.

\section{From intersection to containment, and the price of the bridge}
\label{sec:bridge}

The Bellman cone is not an equivalence class: even after removing scale and
additive symmetries it can contain rewards with genuinely different policy
orderings. Let $F^\circ(\pi):=\{\theta\in K(\pi):\Dnorm{\theta}=1\}$ be the
normalized inverse fiber, modulo the zero directions of $\Dnorm{\cdot}$.

\begin{assumption}[Certified inverse-fiber width]
\label{ass:fiber-width}
There is a known $\eta\in[0,2]$ with
$\sup_{u,v\in F^\circ(\pi)}\Dnorm{u-v}\le\eta$ almost surely for the deployment
population, and each latent reward is behaviorally nonnull.
\end{assumption}

This is the extra information needed to turn intersection into containment. It
must come from side information, multiple environments, or active comparisons;
with none of those, the only universal certificate is $\eta=2$.

\begin{theorem}[Two-stage containment and safe prescription]
\label{thm:containment}
Set $R:=\min\{2,2q_\gamma+\eta\}$ if $q_\gamma<1$ and $R:=2$ otherwise. Then
$\Pr[\Dnorm{\hat\theta_{\mathrm{new}}^\circ-\bar\theta^\circ}\le R]\ge\gamma$,
and if $\mu_R$ maximizes worst-case value over $B_R$ relative to a fixed
reference $\mu_{\mathrm{ref}}\in M(P)$, then
$\Pr[V_{\hat\theta_{\mathrm{new}}}(\mu_R)\ge
V_{\hat\theta_{\mathrm{new}}}(\mu_{\mathrm{ref}})]\ge\gamma$.
\end{theorem}

The factor $2$ is the price of keeping Stage I convex. The reference is not
cosmetic: $B_R$ is \emph{unbounded}---it contains all of
$\bar\theta^\circ+\mathcal A(P)$---so worst-case \emph{absolute} value is
$-\infty$, and only a comparison between two occupancy measures is insensitive
to directions with no behavioral content. Relative to $\mu_{\mathrm{ref}}$ the
inner minimum is the center advantage minus $R\,\gamma_{D_F}(x)$, where
$\gamma_{D_F}$ is the Minkowski gauge polar to $\Dnorm{\cdot}$; it has an LP
epigraph, so the deployed problem is a \emph{single} linear program with penalty
normalized to $[0,1]$ (Appendix~\ref{app:angular}).

\begin{proposition}[Sharp degeneracy threshold]
\label{prop:degeneracy}
For the unit-span center, if $R\ge1$ the optimal relative robust value is zero,
attained by $\mu_{\mathrm{ref}}$; if $R>1$ every optimizer is
feature-equivalent to the reference. A nontrivial safe improvement therefore
requires $2q_\gamma+\eta<1$.
\end{proposition}

This turns vacuity into an auditable output: the implementation reports
$q_\gamma$, $\eta$, $R$, and whether $R<1$. All experiments center
$\mu_{\mathrm{ref}}$ at the calibration population's mean occupancy, which asks
the question \method{} is about---does this policy beat what a random draw from
the population achieves?

\section{Experiments}
\label{sec:experiments}

Both experiments are mechanism tests, not a claim that robustification lowers
true-reward regret. We use random tabular MDPs, a Voronoi-feature gridworld, and
Objectworld ($|S|\le36$, $d\le5$), fitting MaxEnt and Bayesian IRL on 20
training demonstrators, calibrating on 20, and testing on 20 more, over 10
environment seeds (1{,}200 test rewards). Intervals are 95\% Student-$t$ over
the 30 independent environment--seed blocks; details in
Appendix~\ref{app:extra}.

\begin{figure}[t]
\centering
\includegraphics[width=0.86\linewidth]{figures/fig1_metric.pdf}
\caption{The two failures. \textbf{(a)} A random invertible feature
reparameterization of condition number $\kappa$ leaves every pointwise reward and
policy unchanged, yet the inherited angular cap moves; the unit-span score stays
fixed to solver tolerance (shading: mean to worst over 10 seeds).
\textbf{(b)} At $\gamma=0.8$ every calibration event attains marginal coverage,
but containment of the latent reward ray at the same radius does not. Error bars
are 95\% $t$ intervals over 30 environment--seed blocks.}
\label{fig:metric}
\end{figure}

\paragraph{E1: a reward-preserving intervention moves the angular radius.} We
draw one random invertible $A$ at each $\kappa\in\{1,2,5,10,50,200\}$ and
transform $\phi\mapsto A\phi$, $\theta\mapsto A^{-\top}\theta$. Coordinates
change; every scalar reward is preserved to $2.2\times10^{-14}$, so a
behaviorally meaningful radius should not move. At $\kappa=200$ the angular
radius changes by $0.424$ radians ($24.3^\circ$) on average and $1.071$ radians
in the worst seed, while the normalized span radius changes by at most
$2.0\times10^{-8}$ (Figure~\ref{fig:metric}a). This is a falsification test, not
a favorable data distribution: the two problems are mathematically identical.

\paragraph{E2: intersection and containment come apart.} At $\gamma=0.8$, pooled
calibration-event rates are $0.812\pm0.042$ (angular), $0.838\pm0.046$
(value-gap), and $0.819\pm0.036$ (decision-aware); the first and third are
exactly inverse-set intersection. Latent-ray containment at the same radius is
only $0.368\pm0.095$ and $0.438\pm0.081$ respectively, and the
value-gap-to-Euclidean conversion contains $0.004\pm0.005$
(Figure~\ref{fig:metric}b). Conformal ranking behaves exactly as the theory
says; the stronger event that robust safety needs does not follow from
$q_\gamma$ alone. The calibrated radius has pooled mean $0.316$, so by
Proposition~\ref{prop:degeneracy} a certificate would need $\eta<0.37$ on
average to leave anything to decide---and the universal $\eta=2$ leaves nothing
(Appendix~\ref{app:extra}).

\section{What this means}
\label{sec:discussion}

Metric choice is real and measurable (E1), but a conformal intersection radius
is not yet a CIO uncertainty set for the latent reward (E2).
Assumption~\ref{ass:fiber-width} names the missing ingredient and
Proposition~\ref{prop:degeneracy} makes its required strength checkable at run
time; the
universal diameter $2$ proves containment but returns the reference, which is
honest safety with no decision content rather than an algorithmic failure. The
open problem is thus not another IRL point estimator but a mechanism certifying
a small $\eta$, and our exchangeability unit---one demonstrator, one exactly
observed optimal policy, a shared known tabular MDP---is itself demanding.

The pattern generalizes: whenever an inverse problem is partially identified, a
calibrate-then-prescribe pipeline must separately answer in which geometry the
score is measured, how the quantile is computed, and what the calibrated event
implies about the latent parameter. Conformal prediction answers only the
second. CIO can leave the others implicit because one-shot inverse-feasible sets
are thin; IRL is the regime where they are not.

\bibliographystyle{plainnat}
\bibliography{refs}
\appendix

\section{Deferred results: robust geometry and the inherited angular construction}
\label{app:angular}

This appendix states the robust-program geometry referred to in
\S\ref{sec:method} and collects the CIO-inherited pipeline that
\S\ref{sec:method} replaces, since the experiments compare against it.

\begin{theorem}[Concavity under reward-only ambiguity; failure with dynamics ambiguity]
\label{thm:convexity}
\textbf{(a)} For \emph{any} nonempty compact $C\subseteq\mathbb R^d$, if $P$ is
fixed and known then $\max_{\mu\in M(P)}\min_{\theta\in C}\theta^\top\Phi^\top\mu$
is a concave maximization over the fixed polytope $M(P)$, hence a convex outer
problem whenever $C$ admits a tractable linear-minimization oracle.
\textbf{(b)} The guarantee need not survive when ambiguity couples reward and
transition: there is a 2-state, 2-action MDP and a convex one-dimensional
non-product ambiguity set over $(\theta,P)$ whose worst-case value function is
not concave.
\end{theorem}

Part (b) is existential (Appendix~\ref{app:proof-convexity}): rectangular and
other structured cases remain tractable, and failure of concavity is not a
hardness proof.

\begin{proposition}[Resolvent-based score reduction]
\label{prop:resolvent}
Fix $\pi$; let $P_\pi(s,s')=P(s,\pi(s),s')$, let $\Phi_\pi$ be the feature
matrix restricted to $\pi$'s actions, and let $W_\pi:=(I-\beta
P_\pi)^{-1}\Phi_\pi$. Then $V^\pi_\theta=W_\pi\theta$ and
$Q^\pi_\theta(s,a)=\Psi_\pi(s,a)^\top\theta$ with
$\Psi_\pi(s,a):=\phi(s,a)+\beta P(s,a,\cdot)W_\pi$, both linear in $\theta$, so
\[
K(\pi)=\{\theta:\
\Psi_\pi(s,a)^\top\theta\le\Psi_\pi(s,\pi(s))^\top\theta\ \forall s,a\}
\]
exactly. The angular baseline additionally uses
$\hat\Theta_2(\pi)=K(\pi)\cap\{\theta:\|\theta\|_2\le1\}$. Both have $d$
decision variables instead of $d+|S|$ and $|S|(|A|-1)$ nontrivial linear
inequalities, after one resolvent solve per distinct policy.
\end{proposition}

\begin{proposition}[Projective span geometry]
\label{prop:projective}
$d_\circ(\theta,\theta'):=\Dnorm{\theta^\circ-\theta'^\circ}$, with
$\theta^\circ:=\theta/\Dnorm{\theta}$, is a pseudometric with range $[0,2]$,
invariant to independent positive rescaling, to addition of elements of
$\mathcal A(P)$, and to invertible feature reparameterization. For nonnull
rewards $d_\circ(\theta,\theta')=0$ iff $\theta'=c\theta+a$ with $c>0$,
$a\in\mathcal A(P)$; and
$\mathrm{Reg}_{\theta}(\pi_{\theta'})/\Dnorm{\theta}\le d_\circ(\theta,\theta')$.
\end{proposition}

\paragraph{Angular calibration.} For this inherited baseline only, we replace every nonzero point estimate by
its Euclidean normalization and hence use $\|\bar\theta\|_2=1$; the
decision-aware construction instead uses intrinsic unit-span normalization. With the
score $c_k=\max_{\theta\in\hat\Theta_2(\pi_k)}\theta^\top\bar\theta$ of
Proposition~\ref{prop:resolvent}:

\begin{theorem}[Coverage]
\label{thm:coverage}
Under Assumption~\ref{ass:unit}, let $\tau=\lceil(N_{\mathrm{val}}+1)\gamma\rceil$.
If $\tau\le N_{\mathrm{val}}$, let $q_\gamma$ be the $\tau$-th \emph{largest}
score among $c_1,\dots,c_{N_{\mathrm{val}}}$; if $\tau=N_{\mathrm{val}}+1$, set
the sentinel $q_\gamma=-1$. With $\alpha_\gamma:=\arccos(q_\gamma)$,
\[
\mathbb P\big(\hat\Theta_{2,\mathrm{new}}\cap C(\bar\theta,\alpha_\gamma)\ne\emptyset\big)\ge\gamma .
\]
This is an intersection-nonemptiness guarantee, not a containment guarantee: it
does not assert $\theta^*_{\mathrm{new}}\in C(\bar\theta,\alpha_\gamma)$.
\end{theorem}

The proof (Appendix~\ref{app:proof-coverage}) transplants CIO's exchangeability
argument onto these scores. A marginal-coverage simulation against the actual
calibration code (150 trials per $(\gamma,N_{\mathrm{val}})$ cell) confirms
coverage $\ge\gamma$ at every tested combination; tied scores from finite
populations are the likely source of the mild over-coverage beyond the rank
argument's $1/(N_{\mathrm{val}}+1)$ slack. The sentinel returns the whole unit
ball and is valid but uninformative; more generally, because the zero reward is
Bellman-feasible for every policy, any cap with $\alpha\ge\pi/2$ makes the
intersection event automatic, so coverage alone never establishes reward
identification.

\begin{theorem}[Regret bound under a reward-identification bridge]
\label{thm:regret}
Let $\nu(\mu):=\|\Phi^\top\mu\|_2/(1-\beta)\le\bar\nu:=\|\phi\|_{2,\infty}/(1-\beta)$.
Let $\theta^*_{\mathrm{new}}$ be a unit-norm deployment reward, $\mu_{\mathrm{CIRL}}$
maximize the worst-case value over $C=C(\bar\theta,\alpha)$, and $\mu^*$ be
optimal for $\theta^*_{\mathrm{new}}$. Suppose (i) $\alpha\in[0,\pi/2]$ and
(ii) \emph{compatible-reward accuracy}: some
$\theta_1\in\hat\Theta_{2,\mathrm{new}}\cap C$ has
$\|\theta_1-\theta^*_{\mathrm{new}}\|_2\le\eta$. Then
\[
\mathrm{AOG}:=V_{\theta^*_{\mathrm{new}}}(\mu^*)-V_{\theta^*_{\mathrm{new}}}(\mu_{\mathrm{CIRL}})
\le(\eta+2\sin\alpha)\big(\nu(\mu^*)+\nu(\mu_{\mathrm{CIRL}})\big)\le2(\eta+2\sin\alpha)\bar\nu .
\]
\end{theorem}

Theorem~\ref{thm:containment} supplies the analogous bridge in the intrinsic
geometry when the normalized inverse-fiber diameter is certified.

\begin{proposition}[The calibrated robust MDP is a regularized MDP]
\label{prop:regularizer}
Let $\|\bar\theta\|_2=1$, $\alpha\in[0,\pi/2)$. Write $x=a(x)\bar\theta+x_\perp$
with $a(x)=\bar\theta^\top x$ and $b(x)=\|x_\perp\|_2$. Then
\[
\min_{\theta\in C(\bar\theta,\alpha)}\theta^\top x=
\begin{cases}
-\|x\|_2, & \text{if } a(x)\le-\|x\|_2\cos\alpha,\\
a(x)\cos\alpha-b(x)\sin\alpha, & \text{otherwise,}
\end{cases}
\]
so whenever every $\mu\in M(P)$ earns nonnegative estimated reward the robust
policy solves $\max_{\mu\in M(P)}\{\bar\theta^\top\Phi^\top\mu-\tan\alpha\cdot
b(\Phi^\top\mu)\}$.
\end{proposition}

Under the sign condition, then, \method{} is the point-estimate MDP with an
$\ell_2$ penalty at weight $\tan\alpha_\gamma$, fixed by the conformal quantile
rather than tuned: $\alpha_\gamma\to0$ recovers the point estimate and
$\alpha_\gamma\to\pi/2$ refuses any off-direction feature component. Outside the
sign condition the piecewise formula, not the regularizer reading, is the
objective. The same form prices the hedge before deployment: writing
$\mu_{\mathrm{plain}}$ for the point-estimate optimum,
\[
0\le\bar\theta^\top\Phi^\top\mu_{\mathrm{plain}}-\bar\theta^\top\Phi^\top\mu_{\mathrm{CIRL}}
\le\tan\alpha_\gamma\cdot b(\Phi^\top\mu_{\mathrm{plain}}),
\]
computable without ground truth. It bounds the loss in \emph{estimated} reward
only.

\paragraph{The value-gap repair and why it is incomplete.} Replacing the
angular score by the suboptimality gap under $\bar\theta$ itself,
$c_k:=V_{\bar\theta}(\mu^*_{\bar\theta})-V_{\bar\theta}(\mu_{\pi_k})$, needs one
value-iteration solve plus one policy evaluation per demonstrator and no
optimization over $\theta$.

\begin{lemma}[Shaping directions are exactly redundant on $M(P)$]
\label{lem:shaping-redundant}
For any shaping direction $c$ (via potential $\Phi_{\mathrm{pot}}$) and
\textbf{any} $\mu\in M(P)$, $c^\top\Phi^\top\mu=-(1-\beta)\rho_0^\top\Phi_{\mathrm{pot}}$,
a constant independent of $\mu$.
\end{lemma}

\begin{theorem}[Shaping invariance of the value-gap score]
\label{thm:vg-shaping}
Let $\bar\theta'=\bar\theta+c$ for $c\in S_{\mathrm{pot}}$. Then $\bar\theta'$
induces the same optimal occupancy measure as $\bar\theta$, automatically by
Lemma~\ref{lem:shaping-redundant} and with no extra hypothesis, and
$c_k(\bar\theta')=c_k(\bar\theta)$ for every $\pi_k$.
\end{theorem}

Converting this score to a ball $C_{\mathrm{ball}}(\bar\theta,\rho)$ with
$\rho:=q_\gamma/\bar\nu$ uses the Lipschitz constant of
Theorem~\ref{thm:regret}, which is where the construction fails: the bound
licenses gap $\le$ const $\times$ distance and never the reverse, so $q_\gamma$
alone certifies nothing about the set (measured containment $0.003$,
Figure~\ref{fig:metric}b). Its robust program is at least unconditional:

\begin{proposition}[The value-ball cap needs no case split]
\label{prop:ball-regularizer}
For any $\rho\ge0$ and $x$,
$\min_{\|\theta-\bar\theta\|_2\le\rho}\theta^\top x=\bar\theta^\top x-\rho\|x\|_2$,
so $\max_{\mu\in M(P)}[\bar\theta^\top\Phi^\top\mu-\rho\|\Phi^\top\mu\|_2]$ is
the point-estimate MDP with a fixed-weight $\ell_2$ penalty, for every
$\bar\theta$ and $\rho$, with no sign condition.
\end{proposition}

\paragraph{A tie-breaking pitfall in the degenerate separation example.}
A direct transplant of CIO's separation example does not survive inspection. In
a one-state, two-action MDP with $\phi_0=(0,-1)$, $\phi_1=(-u,0)$ and
demonstrator rewards on the first-quadrant arc, the exact-fit LP estimator
returns the tie direction $\theta_u=(1,u)/\sqrt{1+u^2}$, under which both
actions are optimal. Every score is then $c_k=1$, calibration returns
$\alpha_\gamma=0$, and the cap collapses to $\{\theta_u\}$, under which every
action mixture has equal value. The apparent separation is produced by the
solver's tie-breaking rule, not by robustification; a valid separation needs a
strictly positive radius and a unique robust optimizer. The same degeneracy
explains why projective experiments must audit whether a fitted center has
positive occupancy span.

\section{Extended related work}
\label{app:related}

CIO calibrates uncertainty for a latent objective
\citep{lin2024conformal} using split conformal prediction
\citep{vovk2005algorithmic}; a parallel IRL line studies partial identifiability
and feasible reward sets \citep{skalse2023invariance,kim2024unified,
poiani2024suboptimal,lazzati2025robustness,kim2026feasible} under observation
models different from ours. Reward-comparison metrics ask which distance between
rewards is meaningful: EPIC \citep{gleave2021epic} cancels shaping, DARD
\citep{wulfe2022dard} restricts to realizable transitions, and STARC
\citep{skalse2024starc} unifies both. Our occupancy span is the canonical
seminorm behind one STARC construction, so the contribution is not a new reward
metric: it is calibrating Bellman cones \emph{in} that geometry, plus the
explicit fiber-width bridge. Robust MDPs classically recover tractability
through rectangularity \citep{iyengar2005robust,wiesemann2013robust};
\citet{gadot2024solving} solve the fixed-kernel $L_p$-ball reward case. We are
not aware of prior work conformally calibrating an inverse-feasible reward set
from a demonstrator population.

\section{Additional experiments}
\label{app:extra}

\paragraph{Setup detail.} E1--E2 use random tabular MDPs, a Voronoi-feature
gridworld, and Objectworld, each with $|S|\le36$ and $d\le5$. We fit MaxEnt IRL
and Bayesian IRL by Metropolis--Hastings on 20 training demonstrators, calibrate
on 20, and evaluate on 20 fresh demonstrators, for 10 environment seeds: 60
fitted/calibrated runs and 1{,}200 test rewards. Draws within a run share a
center and radius, so intervals are 95\% Student-$t$ over the 30 independent
environment--seed blocks after averaging the two estimators within each block.
Computation used a 20-core laptop, no GPU, and no commercial solver.

\paragraph{The identification budget (E3).} The calibrated $q_\gamma$ has pooled
mean $0.316$, median $0.289$, and range $[0.019,0.871]$, so by
Proposition~\ref{prop:degeneracy} a certified width must satisfy
$\eta<\max\{1-2q_\gamma,0\}$ for the safe prescription to differ from the
reference. With hypothetical certificates $\eta=0$, $0.2$, and $0.5$,
respectively 85.0\%, 76.7\%, and 36.7\% of runs remain nontrivial, and the
environment matters: at $\eta=0.2$ the fractions are 95\% for random MDPs, 60\%
for gridworld, and 75\% for Objectworld. The universal certificate $\eta=2$
instead gives $R=2$ in every run, expanded containment $1.000$, reference
dominance $1.000$, and nontriviality $0.000$---a theorem/implementation check,
not a useful policy, since the deployed occupancy is exactly the reference
(Figure~\ref{fig:identification-frontier}).

\begin{figure}[t]
\centering
\includegraphics[width=0.92\linewidth]{figures/fig2_identification_frontier.pdf}
\caption{Identification-budget frontier for the numbers quoted in
\S\ref{sec:experiments}. \textbf{(a)} Distribution of the dimensionless Stage-I
radius over 20 runs per environment. Even a zero-width fiber can be nontrivial
only below $q_\gamma=0.5$. \textbf{(b)} For each hypothetical certified width
$\eta$, the fraction of runs satisfying $2q_\gamma+\eta<1$. The curve does not
estimate or validate $\eta$; it states how strong an independent identification
mechanism must be. Shading is a 95\% $t$ interval over 30 environment--seed
blocks.}
\label{fig:identification-frontier}
\end{figure}

\paragraph{Projective-center eligibility.} We refit the exact-fit
suboptimality-loss LP on the same 30 environment--seed training populations as
the main audit. Its center has positive occupancy span in 29 cases and is
additively null in one gridworld case. The latter has no unit-span representative.
The MaxEnt and Bayesian estimators used in E2--E3 are nonnull in every reported
run. This check is a domain-of-definition audit, not a comparison of estimator
accuracy.

\paragraph{Legacy policy-regret comparison.} Before introducing the certified
fiber-width bridge, the exploratory three-estimator audit deployed each method
at its Stage-I radius. Angular robustification increased true-reward regret in
all nine environment--estimator cells, by factors $1.12$--$19.49$. The
value-gap ball was $0.72$--$1.86\times$ the point-estimate regret and improved
three cells; the decision-aware radius improved five. These policies are not
the Stage-II policy of Theorem~\ref{thm:containment}, so the table diagnoses
geometry and objective mismatch rather than demonstrating safe improvement.

\paragraph{Concentration baseline.} Exact conformal calibration is tighter than
a DKW bound everywhere tested; on the unbounded value-gap score DKW inflates
true-reward regret $4$--$460\times$ over the point estimate, because the same
additive slack $\sqrt{\log(2/\delta)/(2N)}$ shifts a raw-value threshold far
more than a cosine in $[-1,1]$.

\paragraph{Scale.} Rerunning the full gridworld pipeline at $|S|=225$, $d=7$
($6.25\times$ more states, two more reward dimensions) leaves the angular
baseline's negative result unchanged: its robust regret exceeds the plain
estimate's for all three estimators (LP-IRL $1.03\to3.19$, MaxEnt
$1.01\to2.88$, Bayesian
$0.06\to0.36$).

\paragraph{RLHF relabeling.} On a ``response selection'' MDP whose population
is annotators rather than task demonstrators, with the calibration and
robust-MDP code unchanged, coverage meets or exceeds target across
$\gamma\in\{0.5,0.7,0.8,0.9\}$ (0.947--0.994, 8 seeds, one configuration).
Across four annotator-population configurations (4--10 responses, 3--6
attributes, 15--60 annotators) and 30 seeds at $\gamma\in\{0.7,0.8,0.9\}$, mean
coverage per cell is $0.966$--$0.991$ with at least 93\% of seeds per cell
meeting target, consistent with the marginal (not per-seed) nature of the
guarantee. The gap above target likely comes from the same score-atom
over-coverage as Theorem~\ref{thm:coverage}, sharpened by small populations.

\paragraph{Beyond linear rewards.} $\Dnorm{\cdot}$ and $s_k$ are instances of an
integral probability metric $\sup_{f\in\mathcal F_\psi}[\mathbb
E_{\mu^*_{\bar\theta}}(f)-\mathbb E_{\mu_k}(f)]$, with
$\mathcal F_\psi=\{\bar\theta\}$ giving the value-gap score,
$\mathcal F_\psi=\{$linear, $\|\theta\|_2\le1\}$ a closed-form dual-norm score,
and richer classes (an RKHS ball, a bounded neural discriminator as in
GAIL/AIRL) requiring adversarial training. On exact (non-sampled) occupancy
measures, projected gradient ascent recovers the linear closed form to solver
precision, and a strictly richer critic class, which contains the linear one and
therefore weakly dominates it, finds $3$--$4\times$ more separation in every
tested case. Whether a finite-trajectory \emph{estimate} concentrates
tightly enough, as $\mathcal F_\psi$ grows, for calibration to remain
non-vacuous is the open piece with no directly reusable existing theory.

\paragraph{Reproduction.} From the repository root, the main artifacts are
generated by
\texttt{experiments/run\_reparam\_invariance.py},
\texttt{experiments/run\_containment\_audit.py}, and
\texttt{experiments/run\_center\_validity.py}; the two figures are generated by
\texttt{paper/make\_fig1\_metric.py} and
\texttt{paper/make\_fig2\_identification\_frontier.py}. All scripts set NumPy
random generators from the recorded integer seed. The audit JSON stores the
schema version, split sizes, target $\gamma$, environment, estimator, and seed
in every row.

\section{Proofs}
\label{app:proofs}

Every identity and inequality stated in the main text is also checked
numerically, independently of the library routines it describes, by
\texttt{paper/verify\_formulas.py}: the two-solve form of $\Dnorm{\cdot}$
against brute-force enumeration of $M(P)$'s vertices, the seminorm axioms, the
regret bound, reparametrization invariance, the zero set $\mathcal A(P)$,
Lemma~\ref{lem:shaping-redundant}, the piecewise cap minimum of
Proposition~\ref{prop:regularizer}, the cap diameter and Lipschitz constant of
Theorem~\ref{thm:regret}, and the gauge form of the deployed robust objective.

\subsection{Proof of Proposition~\ref{prop:resolvent}}
\label{app:proof-resolvent}

The value function $V^\pi_\theta$ solves the linear Bellman equation:
\[
V^\pi_\theta = \Phi_\pi\theta + \beta P_\pi V^\pi_\theta.
\]
Since $\beta\in[0,1)$ and $P_\pi$ is a stochastic matrix, its spectral radius satisfies $\rho(\beta P_\pi)\le\beta<1$, so $I-\beta P_\pi$ is invertible via a Neumann series, yielding:
\[
V^\pi_\theta = (I-\beta P_\pi)^{-1}\Phi_\pi\theta = W_\pi\theta.
\]
Substituting $V^\pi_\theta$ into the Bellman equation for the action-value function $Q^\pi_\theta(s,a)$ gives:
\begin{align*}
Q^\pi_\theta(s,a) &= \phi(s,a)^\top\theta + \beta P(s,a,\cdot) V^\pi_\theta \\
&= \phi(s,a)^\top\theta + \beta P(s,a,\cdot) W_\pi\theta \\
&= \Psi_\pi(s,a)^\top\theta,
\end{align*}
where $\Psi_\pi(s,a) := \phi(s,a) + \beta P(s,a,\cdot) W_\pi$, which is linear in $\theta$ as claimed.

By the policy improvement theorem, $\pi$ is Bellman-optimal for reward $\theta$ if and only if
\[
Q^\pi_\theta(s,a) \leq Q^\pi_\theta(s,\pi(s)) \quad \forall s \in S, \forall a \in A.
\]
Substituting the linear expression for $Q^\pi_\theta(s,a)$ yields exactly the stated cone:
\[
K(\pi) = \{\theta : \Psi_\pi(s,a)^\top\theta \leq \Psi_\pi(s,\pi(s))^\top\theta\ \ \forall s,a\},
\]
with no relaxation since the substitution is an equality.

The calibration score
\[
c_k = \max_{\theta\in\hat\Theta_2(\pi_k)}\theta^\top\bar\theta
\]
is then a linear objective over a set of $|S||A|$ linear inequalities plus the unit-ball constraint in $\theta\in\mathbb{R}^d$ alone: an SOCP of size $d$, independent of $|S|$ once $\Psi_{\pi_k}$ is precomputed. $\blacksquare$

\subsection{Proof of Theorem~\ref{thm:convexity}}
\label{app:proof-convexity}

\textbf{(a)} Define $g_C(x) := \min_{\theta\in C}\theta^\top x$. This is a pointwise infimum of the affine family $\{x\mapsto\theta^\top x : \theta\in C\}$. Since a pointwise infimum of affine functions is concave by definition for any nonempty $C$ (convex, nonconvex, or discrete; concavity of $g_C$ requires no assumption on $C$'s shape), composing $g_C$ with the linear map $\mu\mapsto\Phi^\top\mu$ preserves concavity:
\[
\mu \mapsto g_C(\Phi^\top\mu) = \min_{\theta\in C}\theta^\top\Phi^\top\mu
\]
is concave in $\mu$.

Since $M(P)$ is a fixed polytope that does not depend on $\theta$ (the transition kernel $P$ is known and fixed), the ambiguity set never touches the feasible region, only the objective. Therefore,
\[
\max_{\mu\in M(P)} g_C(\Phi^\top\mu)
\]
is a concave maximization over a fixed compact convex set for \emph{any}
nonempty compact $C$. This establishes the convex geometry; efficient solution
also requires an efficient representation or linear-minimization oracle for
$C$. No Bellman recursion or rectangularity assumption enters the geometric
argument.

\textbf{(b)} Let state $s_0$ be active and $s_1$ absorbing with zero reward.
At $s_0$, action $a_0$ earns zero and moves to $s_1$; action $a_1$ earns
$1+t$ and returns to $s_0$ with probability $1-t$ (otherwise moving to
$s_1$), where the same ambiguity parameter $t\in[0,1]$ couples reward and
dynamics. Both actions at $s_1$ self-loop with zero reward. Starting from
$s_0$, let a stationary policy choose $a_1$ with probability $p$. Its
worst-case discounted value is
\[
W(p) = \min_{t\in[0,1]} \frac{p(1+t)}{1-\beta p\,q(t)} = \min_{t\in[0,1]} \frac{p(1+t)}{1-\beta p(1-t)}.
\]
At discount factor $\beta=0.9$, evaluating $W(p)$ at three points gives:
\begin{align*}
W(0) &= 0, \\
W(5/9) &= \frac{10}{9}, \quad \text{constant in }t, \\
W(5/18) &= \min_{t\in[0,1]}
  \frac{10(1+t)}{27+9t} = \frac{10}{27}
  \quad (\text{attained at }t=0).
\end{align*}
Since $5/18 = \tfrac12(0 + 5/9)$, concavity of $W$ would require:
\[
W(5/18) \geq \frac{1}{2}\big[W(0) + W(5/9)\big] = \frac{5}{9}.
\]
But $10/27<5/9$, producing a strict rational-valued Jensen violation.

Hence $W$ is not concave, and the analogous worst-case value function for this jointly-ambiguous $(\theta,P)$ instance is not concave either, so the outer maximization is not a convex program in this regime. $\blacksquare$

\subsection{Proof of Proposition~\ref{prop:da-seminorm}}
\label{app:proof-da-seminorm}

\textbf{Seminorm properties.} Nonnegativity is immediate ($\max\ge\min$).
Positive homogeneity: $V_v(\mu)$ is linear in $v$, so for $c>0$,
$\max_\mu V_{cv}(\mu)-\min_\mu V_{cv}(\mu) = c\big(\max_\mu
V_v(\mu)-\min_\mu V_v(\mu)\big)=c\Dnorm{v}$. Subadditivity: since
$\sup$ of a sum is at most the sum of $\sup$s and $\inf$ of a sum is at least
the sum of $\inf$s,
$\Dnorm{u+v} = \max_\mu V_{u+v}(\mu)-\min_\mu V_{u+v}(\mu) \le
\big(\max_\mu V_u(\mu)+\max_\mu V_v(\mu)\big) - \big(\min_\mu
V_u(\mu)+\min_\mu V_v(\mu)\big) = \Dnorm{u}+\Dnorm{v}$.

\textbf{Zero set.} $\Dnorm{v}=0$ iff $V_v(\mu)$ is constant over
$M(P)$ iff $v^\top\Phi^\top(\mu-\nu)=0$ for every $\mu,\nu\in M(P)$, which is
exactly the definition of $v\in\mathcal A(P)$. Every shaping direction
$c\in S_{\mathrm{pot}}$ satisfies $c^\top\Phi^\top(\mu-\nu)=0$ for all
$\mu,\nu\in M(P)$ by Lemma~\ref{lem:shaping-redundant} (the constant
$-(1-\beta)\rho_0^\top\Phi_{\mathrm{pot}}$ cancels in the difference), so
$S_{\mathrm{pot}}\subseteq\mathcal A(P)$. The containment can be strict: if a
state $s$ is unreachable under every policy then $\mu(s,\cdot)=0$ throughout
$M(P)$, so any $v$ whose induced reward is supported on $s$ lies in $\mathcal
A(P)$, while such a reward need not be a potential-shaping term. We do not
claim strictness in general, only that $\mathcal A(P)$ is the exact zero set of
$\Dnorm{\cdot}$, which is what the invariance claims use.

\textbf{Regret bound.} With $\mu^*_\theta$ optimal for $\theta$ and
$\mu_{\theta'}$ optimal for $\theta'$,
\begin{align*}
\mathrm{Reg}_\theta(\pi_{\theta'})
&= V_\theta(\mu^*_\theta) - V_\theta(\mu_{\theta'}) \\
&= \big[V_\theta(\mu^*_\theta)-V_{\theta'}(\mu^*_\theta)\big]
  - \big[V_\theta(\mu_{\theta'})-V_{\theta'}(\mu_{\theta'})\big]
  + \big[V_{\theta'}(\mu^*_\theta)-V_{\theta'}(\mu_{\theta'})\big] \\
&= V_{\theta-\theta'}(\mu^*_\theta) - V_{\theta-\theta'}(\mu_{\theta'})
  + \underbrace{\big[V_{\theta'}(\mu^*_\theta)-V_{\theta'}(\mu_{\theta'})\big]}_{\le\,0
  \text{ since $\mu_{\theta'}$ maximizes }V_{\theta'}\text{ over }M(P)} \\
&\le V_{\theta-\theta'}(\mu^*_\theta) - V_{\theta-\theta'}(\mu_{\theta'})
\;\le\; \max_\mu V_{\theta-\theta'}(\mu) - \min_\mu V_{\theta-\theta'}(\mu)
\;=\; \Dnorm{\theta-\theta'},
\end{align*}
the last step because $\mu^*_\theta,\mu_{\theta'}\in M(P)$, so their
$V_{\theta-\theta'}$ values each lie in $[\min_\mu V_{\theta-\theta'}(\mu),
\max_\mu V_{\theta-\theta'}(\mu)]$.

\textbf{Reparametrization invariance.} For invertible $A$, writing
$\tilde\phi=A\phi$, $\tilde v = A^{-\top}v$: $\tilde\phi(s,a)^\top\tilde v =
\phi(s,a)^\top A^\top A^{-\top}v = \phi(s,a)^\top v$ for every $(s,a)$, so the
reward function is unchanged, hence $V_v(\mu)$ at every $\mu\in M(P)$, hence
$\Dnorm{v}$ itself. $\blacksquare$

\subsection{Proof of Proposition~\ref{prop:projective}}

On the quotient by $\mathcal A(P)$, $\Dnorm{\cdot}$ is a norm and every
nonnull positive ray has a unit representative. Hence the triangle inequality,
symmetry, and nonnegativity of $d_\circ$ follow from the seminorm, and
$d_\circ\le\Dnorm{\theta^\circ}+\Dnorm{\theta'^\circ}=2$. Independent positive
rescaling does not change either unit representative. Adding an element of
$\mathcal A(P)$ changes a representative only by a seminorm-zero vector, and
feature invariance follows from Proposition~\ref{prop:da-seminorm}.

If $d_\circ(\theta,\theta')=0$, then
$\theta/\Dnorm{\theta}-\theta'/\Dnorm{\theta'}\in\mathcal A(P)$; rearranging
gives $\theta'=c\theta+a$ with
$c=\Dnorm{\theta'}/\Dnorm{\theta}>0$ and $a\in\mathcal A(P)$. The converse is
immediate. Finally apply Proposition~\ref{prop:da-seminorm} to the two normalized
rewards. Their optimal policies are unchanged by positive normalization, so
\[
 \mathrm{Reg}_{\theta}(\pi_{\theta'})/\Dnorm{\theta}
 =\mathrm{Reg}_{\theta^\circ}(\pi_{\theta'^\circ})
 \le\Dnorm{\theta^\circ-\theta'^\circ}=d_\circ(\theta,\theta').
 \qquad\blacksquare
\]

\subsection{Proof of Theorem~\ref{thm:da-coverage}}
\label{app:proof-da-coverage}

\textbf{Equivalence.} By definition $s_k=\min_{\theta\in K(\pi_k)}
\Dnorm{\theta-\bar\theta^\circ}$, so $s_k\le q$ iff some
$\theta\in K(\pi_k)$ achieves $\Dnorm{\theta-\bar\theta^\circ}\le
q$, i.e.\ iff $K(\pi_k)\cap B_q\ne\emptyset$. This is an
immediate consequence of $s_k$'s definition as a minimum over the cone, not a
separate argument. The epigraph formulation is a feasible, bounded-below LP
over a closed polyhedron, so the fundamental theorem of linear programming
gives an attained optimum. The implementation adds a large numerical safety
cap and rejects a solution if that cap is active.

\textbf{Coverage.} $\bar\theta$ is a fixed, measurable function of the
training split, disjoint from the calibration split and the fresh
demonstrator. Conditional on $\bar\theta$, each $s_k$ is a deterministic,
measurable function of one i.i.d.\ demonstrator draw, so
$s_1,\dots,s_{N_{\mathrm{val}}},s_{\mathrm{new}}$ are exchangeable. The
ascending rank of $s_{\mathrm{new}}$ among these $N_{\mathrm{val}}+1$
exchangeable reals is therefore uniform on $\{1,\dots,N_{\mathrm{val}}+1\}$,
so $\mathbb P[\mathrm{rank}(s_{\mathrm{new}})\le\tau] =
\tau/(N_{\mathrm{val}}+1)\ge\gamma$ by the choice of $\tau$; $s_{\mathrm{new}}
\le q_\gamma$ iff $\mathrm{rank}(s_{\mathrm{new}})\le\tau$ (ties broken
uniformly at random, as in the proof of Theorem~\ref{thm:coverage}).
$\blacksquare$

\subsection{Proof of Theorem~\ref{thm:containment}}
\label{app:proof-containment}

On the Stage-I event, choose a witness
$\tilde\theta\in K(\pi_{\mathrm{new}})$ with
$\Dnorm{\tilde\theta-\bar\theta^\circ}\le q_\gamma$. If $q_\gamma<1$, the
reverse triangle inequality gives
$\Dnorm{\tilde\theta}\ge1-q_\gamma>0$. Since $K(\pi)$ is a cone,
$\tilde\theta^\circ=\tilde\theta/\Dnorm{\tilde\theta}$ belongs to
$F^\circ(\pi_{\mathrm{new}})$. Moreover,
\begin{align*}
\Dnorm{\tilde\theta^\circ-\bar\theta^\circ}
&\le \Dnorm{\tilde\theta^\circ-\tilde\theta}
   +\Dnorm{\tilde\theta-\bar\theta^\circ}\\
&=|1-\Dnorm{\tilde\theta}|+q_\gamma\le2q_\gamma.
\end{align*}
The latent normalized reward is in the same fiber, so Assumption
\ref{ass:fiber-width} and the triangle inequality give
$\Dnorm{\hat\theta_{\mathrm{new}}^\circ-\bar\theta^\circ}
\le\eta+2q_\gamma$. If $q_\gamma\ge1$, the universal projective diameter $2$
gives the stated fallback. Theorem~\ref{thm:da-coverage} assigns probability at
least $\gamma$ to the Stage-I event, proving containment.

For the robust statement, $\mu_{\mathrm{ref}}$ is feasible and has relative
value zero under every reward, so the optimum at $\mu_R$ is nonnegative. On the
containment event, $\hat\theta_{\mathrm{new}}^\circ\in B_R$; hence its relative
value at $\mu_R$ is at least the worst-case relative value and therefore at
least zero. Positive rescaling back to $\hat\theta_{\mathrm{new}}$ preserves
the inequality. Subtracting both values from the same optimal value gives the
equivalent regret statement. $\blacksquare$

\subsection{Proof of Proposition~\ref{prop:degeneracy}}

Let $x=\Phi^\top(\mu-\mu_{\mathrm{ref}})/(1-\beta)$. Polar duality and
$\Dnorm{\bar\theta^\circ}=1$ give
$V_{\bar\theta^\circ}(\mu)-V_{\bar\theta^\circ}(\mu_{\mathrm{ref}})
\le\gamma_{D_F}(x)$. Therefore the relative robust objective is at most
$(1-R)\gamma_{D_F}(x)\le0$ when $R\ge1$. The reference has $x=0$ and attains
zero. If $R>1$, equality requires $\gamma_{D_F}(x)=0$, which means the two
occupancies have the same feature vector modulo the null directions; for
$R=1$, additional exposed-face ties can attain zero. $\blacksquare$

\subsection{Proof of Theorem~\ref{thm:coverage}}
\label{app:proof-coverage}

Under Assumption~\ref{ass:unit}, the sequence of pairs
\[
(\hat\theta_1,\pi_1),\,\dots,\,(\hat\theta_{N_{\mathrm{val}}},\pi_{N_{\mathrm{val}}}),\,(\hat\theta_{\mathrm{new}},\pi_{\mathrm{new}})
\]
are i.i.d., and therefore exchangeable. The calibration score
\[
c_k = \max_{\theta\in\hat\Theta_2(\pi_k)}\theta^\top\bar\theta
\]
(Proposition~\ref{prop:resolvent}) is a fixed measurable function of $(\hat\theta_k,\pi_k)$ once $\bar\theta$ is fixed (computed on a disjoint training split, independent of the calibration and test data). Consequently, the calibration scores
\[
c_1,\,\dots,\,c_{N_{\mathrm{val}}},\,c_{\mathrm{new}}
\]
are themselves exchangeable.

Break ties at random only for the rank argument (the implemented inclusive
inequality remains conservative under ties). Exchangeability makes the
descending rank of $c_{\mathrm{new}}$ uniform on
$\{1,\dots,N_{\mathrm{val}}+1\}$. If $\tau\le N_{\mathrm{val}}$, this gives
\[
\mathbb{P}(c_{\mathrm{new}} \geq q_\gamma)
\geq \frac{\tau}{N_{\mathrm{val}}+1} \geq \gamma.
\]
If $\tau=N_{\mathrm{val}}+1$, then $q_\gamma=-1$ and the event holds surely
because $\|\bar\theta\|_2=1$ and every score lies in $[0,1]$ (the zero reward
is feasible for every policy).

It remains to show that the event
$\{c_{\mathrm{new}} \geq q_\gamma\} =
\{c_{\mathrm{new}} \geq \cos\alpha_\gamma\}$ is equivalent to
$\hat\Theta_{2,\mathrm{new}} \cap C(\bar\theta,\alpha_\gamma) \neq \emptyset$:
\begin{align*}
c_{\mathrm{new}} \geq \cos\alpha_\gamma &\iff \exists \theta \in \hat\Theta_{2,\mathrm{new}} \text{ such that } \theta^\top\bar\theta \geq \cos\alpha_\gamma \\
&\iff \exists \theta \in \hat\Theta_{2,\mathrm{new}} \text{ such that } \theta \in C(\bar\theta,\alpha_\gamma) \\
&\iff \hat\Theta_{2,\mathrm{new}} \cap C(\bar\theta,\alpha_\gamma) \neq \emptyset.
\end{align*}
Combining these statements yields:
\[
\mathbb{P}\big(\hat\Theta_{2,\mathrm{new}} \cap C(\bar\theta,\alpha_\gamma) \neq \emptyset\big) \geq \gamma. \quad \blacksquare
\]

\subsection{Proof of Theorem~\ref{thm:regret}}
\label{app:proof-regret}

Write
\[
V_\theta(\mu) := \frac{\theta^\top\Phi^\top\mu}{1-\beta}
\]
for the expected discounted return of occupancy measure $\mu$ under reward $\theta$. By Cauchy--Schwarz, this function is $\bar\nu$-Lipschitz in $\theta$ (in $\ell_2$) uniformly over $\mu\in M(P)$:
\begin{align*}
|V_\theta(\mu) - V_{\theta'}(\mu)| &= \frac{|(\theta-\theta')^\top \Phi^\top\mu|}{1-\beta} \\
&\leq \|\theta-\theta'\|_2 \frac{\|\Phi^\top\mu\|_2}{1-\beta} \\
&\leq \|\theta-\theta'\|_2 \,\nu(\mu).
\end{align*}

Let $\theta^-_*\in\arg\min_{\theta\in C}V_\theta(\mu^*)$ and
$\theta^-_{\mathrm{CIRL}}\in\arg\min_{\theta\in C}
V_\theta(\mu_{\mathrm{CIRL}})$; compactness of $C$ ensures both exist. Robust
optimality gives the correctly quantified comparison
\[
V_{\theta^-_{\mathrm{CIRL}}}(\mu_{\mathrm{CIRL}})
\ge V_{\theta^-_*}(\mu^*).
\]

For $\alpha\le\pi/2$, the Euclidean diameter of the unit-ball cap is at most
$2\sin\alpha$: its extreme separation occurs between two unit vectors on the
rim whose angular separation is $2\alpha$. Because $\theta_1$, $\theta^-_*$,
and $\theta^-_{\mathrm{CIRL}}$ all lie in $C$, Assumption (ii) and the triangle
inequality imply
\[
\|\theta^-_*-\theta^*_{\mathrm{new}}\|_2,
\ \|\theta^-_{\mathrm{CIRL}}-\theta^*_{\mathrm{new}}\|_2
\le 2\sin\alpha+\eta.
\]
We may therefore decompose
\begin{align*}
\mathrm{AOG}
&=V_{\theta^*_{\mathrm{new}}}(\mu^*)
  -V_{\theta^*_{\mathrm{new}}}(\mu_{\mathrm{CIRL}})\\
&\le
  \big[V_{\theta^*_{\mathrm{new}}}(\mu^*)-V_{\theta^-_*}(\mu^*)\big]
  +\big[V_{\theta^-_{\mathrm{CIRL}}}(\mu_{\mathrm{CIRL}})
       -V_{\theta^*_{\mathrm{new}}}(\mu_{\mathrm{CIRL}})\big]\\
&\le (\eta+2\sin\alpha)
  \big(\nu(\mu^*)+\nu(\mu_{\mathrm{CIRL}})\big).
\end{align*}
The uniform corollary follows by replacing both local Lipschitz constants with
$\bar\nu$. $\blacksquare$

\subsection{Proof of Proposition~\ref{prop:regularizer}}
\label{app:proof-regularizer}

If $\alpha=0$, the cap is the singleton $\{\bar\theta\}$ and both claims are
immediate. Assume henceforth that $0<\alpha<\pi/2$.

Fix $x \in \mathbb{R}^d$ and write $a = \bar\theta^\top x$, $x_\perp = x -
a\bar\theta$, $b = \|x_\perp\|_2$, so $\|x\|_2^2 = a^2+b^2$. Any
$\theta$ with $\|\theta\|_2\le1$ decomposes as $\theta = u\bar\theta + v\hat
w$ with $\hat w \perp \bar\theta$ a unit vector, $u = \bar\theta^\top\theta$,
$v \geq 0$, and $u^2+v^2 = \|\theta\|_2^2 \le 1$. The cap constraint
$\theta^\top\bar\theta \ge \cos\alpha$ reads $u \ge \cos\alpha$. Since
$\theta^\top x = ua + v\,\hat w^\top x_\perp$ and $\hat w$ enters only through
$\hat w^\top x_\perp \in [-b, b]$, the minimizing choice is $\hat w =
-x_\perp/b$ when $b>0$ (any $\hat w$ when $b=0$), giving $\theta^\top x = au -
bv$. The inner problem is therefore the two-variable program
\[
\min\ \{\,au - bv \;:\; u \ge \cos\alpha,\ v\ge0,\ u^2+v^2\le1\,\}.
\]
As $b \ge 0$, it is optimal to take $v$ maximal for the chosen $u$, i.e.\
$v = \sqrt{1-u^2}$, leaving $\min_{u\in[\cos\alpha,1]} f(u)$ with $f(u) := au -
b\sqrt{1-u^2}$ and $f'(u) = a + bu/\sqrt{1-u^2}$, which is strictly increasing
in $u$ on $(-1,1)$. Hence $f$ is either increasing on the whole interval (the
constraint $u \ge \cos\alpha$ binds) or has a unique interior stationary point.

\emph{Case 1: $f'(\cos\alpha) \ge 0$}, i.e.\ $a + b\cot\alpha \ge 0$. Then $f$
is nondecreasing on $[\cos\alpha,1]$ and the minimum is at $u = \cos\alpha$,
$v = \sin\alpha$, giving $f = a\cos\alpha - b\sin\alpha$.

\emph{Case 2: $f'(\cos\alpha) < 0$.} Setting $f'(u^\star)=0$ gives
$u^\star/\sqrt{1-{u^\star}^2} = -a/b =: c > 0$, so $u^\star = c/\sqrt{1+c^2}$
and $f(u^\star) = (ac-b)/\sqrt{1+c^2} = -\sqrt{a^2+b^2} = -\|x\|_2$, the
unconstrained minimum over the unit ball, attained at $\theta = -x/\|x\|_2$.
This is feasible exactly when $-x/\|x\|_2$ lies in the cap, i.e.\ $a \le
-\|x\|_2\cos\alpha$, which is equivalent to $f'(\cos\alpha)<0$ after squaring.
The two cases agree at the boundary $a = -\|x\|_2\cos\alpha$, where $b =
\|x\|_2\sin\alpha$ and $a\cos\alpha - b\sin\alpha = -\|x\|_2(\cos^2\alpha +
\sin^2\alpha) = -\|x\|_2$, so the stated piecewise formula is continuous.

For the second claim, every nonzero $x=\Phi^\top\mu$ with $a(x)\ge0$ lies in
Case 1, while at $x=0$ both displayed branches equal zero. Hence the robust objective is
$\mu\mapsto a\cos\alpha - b\sin\alpha$ uniformly on $M(P)$. Since $\alpha <
\pi/2$ gives $\cos\alpha>0$, dividing the objective by the positive constant
$\cos\alpha$ leaves the $\arg\max$ unchanged and yields $a - \tan\alpha\, b$,
which is the stated program. Concavity is immediate from this form:
$a(\Phi^\top\mu) = \bar\theta^\top\Phi^\top\mu$ is linear in $\mu$, while
$b(\Phi^\top\mu)$ is a seminorm composed with a linear map, hence convex, and
$\tan\alpha \ge 0$. $\blacksquare$

\subsection{Proof of Lemma~\ref{lem:shaping-redundant}}
\label{app:proof-shaping-redundant}

Expand, using the definition of a shaping direction,
\begin{align*}
c^\top\Phi^\top\mu &= \sum_{s,a}\mu(s,a)\big[\beta P(s,a,\cdot)^\top\Phi_{\mathrm{pot}}
- \Phi_{\mathrm{pot}}(s)\big] \\
&= \sum_{s'}\Phi_{\mathrm{pot}}(s')\Big[\beta\sum_{s,a}P(s,a,s')\mu(s,a)\Big]
- \sum_s\Phi_{\mathrm{pot}}(s)\sum_a\mu(s,a).
\end{align*}
The flow-conservation constraint defining $M(P)$ gives, for every $s'$,
$\sum_a\mu(s',a) - \beta\sum_{s,a}P(s,a,s')\mu(s,a) = (1-\beta)\rho_0(s')$, i.e.
$\beta\sum_{s,a}P(s,a,s')\mu(s,a) = \sum_a\mu(s',a) - (1-\beta)\rho_0(s')$.
Substituting into the first term,
\[
c^\top\Phi^\top\mu = \underbrace{\sum_{s'}\Phi_{\mathrm{pot}}(s')\sum_a\mu(s',a)}_{=\sum_s\Phi_{\mathrm{pot}}(s)\sum_a\mu(s,a)}
- (1-\beta)\sum_{s'}\Phi_{\mathrm{pot}}(s')\rho_0(s')
- \sum_s\Phi_{\mathrm{pot}}(s)\sum_a\mu(s,a),
\]
and the first and third terms are identical, hence cancel, leaving
$c^\top\Phi^\top\mu = -(1-\beta)\,\rho_0^\top\Phi_{\mathrm{pot}}$, independent of
$\mu$. This uses only the linear flow-conservation constraints defining
$M(P)$, so it holds at \emph{every} $\mu\in M(P)$, not only at policy-occupancy
vertices. $\blacksquare$

\subsection{Proof of Theorem~\ref{thm:vg-shaping}}
\label{app:proof-vg-shaping}

By Lemma~\ref{lem:shaping-redundant}, $\theta'^\top\Phi^\top\mu =
\bar\theta^\top\Phi^\top\mu + c^\top\Phi^\top\mu = \bar\theta^\top\Phi^\top\mu -
(1-\beta)\rho_0^\top\Phi_{\mathrm{pot}}$ for \emph{every} $\mu\in M(P)$, so the
two objectives $\mu\mapsto\bar\theta^\top\Phi^\top\mu$ and
$\mu\mapsto\bar\theta'^\top\Phi^\top\mu$ differ by a $\mu$-independent additive
constant. Hence they share the same maximizer, $\mu^*_{\bar\theta'} =
\mu^*_{\bar\theta}$, and
$V_{\bar\theta'}(\mu^*_{\bar\theta'}) = V_{\bar\theta}(\mu^*_{\bar\theta}) -
\rho_0^\top\Phi_{\mathrm{pot}}$. The same additive-constant identity applied at
$\mu=\mu_{\pi_k}$ gives $V_{\bar\theta'}(\mu_{\pi_k}) =
V_{\bar\theta}(\mu_{\pi_k}) - \rho_0^\top\Phi_{\mathrm{pot}}$. Subtracting,
\[
c_k(\bar\theta') = \big[V_{\bar\theta}(\mu^*_{\bar\theta}) -
\rho_0^\top\Phi_{\mathrm{pot}}\big] - \big[V_{\bar\theta}(\mu_{\pi_k}) -
\rho_0^\top\Phi_{\mathrm{pot}}\big] = c_k(\bar\theta). \quad\blacksquare
\]

\subsection{Proof of Proposition~\ref{prop:ball-regularizer}}
\label{app:proof-ball-regularizer}

Write $\theta=\bar\theta+\delta$, $\|\delta\|_2\le\rho$. Then
$\theta^\top x = \bar\theta^\top x + \delta^\top x$, and
$\min_{\|\delta\|_2\le\rho}\delta^\top x = -\rho\|x\|_2$ by Cauchy--Schwarz
(equality at $\delta=-\rho x/\|x\|_2$ when $x\neq0$, and trivially at $x=0$),
giving $\min_{\theta:\|\theta-\bar\theta\|_2\le\rho}\theta^\top x =
\bar\theta^\top x - \rho\|x\|_2$ for every $x$, with no case split. Substituting
$x=\Phi^\top\mu$ gives the stated robust MDP objective directly; both terms
are defined for every $\bar\theta,\rho\ge0,\mu\in M(P)$, so no sign condition
on $\bar\theta^\top\Phi^\top\mu$ is needed (contrast
Proposition~\ref{prop:regularizer}, whose regularizer form requires the
nonnegative-return branch). $\blacksquare$

\end{document}
